\documentclass[11pt,a4paper]{article}
\usepackage[utf8]{inputenc}
\usepackage[T1]{fontenc}
\usepackage{geometry}
\geometry{margin=1in}
\usepackage{hyperref}
\usepackage{enumitem}
\usepackage{booktabs}
\usepackage{longtable}
\usepackage{xcolor}
\usepackage{titlesec}
\usepackage{amsmath,amssymb}

\titleformat{\section}{\Large\bfseries\color{blue!70!black}}{}{0em}{}
\titleformat{\subsection}{\large\bfseries\color{blue!50!black}}{}{0em}{}
\titleformat{\subsubsection}{\normalsize\bfseries\color{blue!30!black}}{}{0em}{}

\title{\textbf{Replication Plan}\\[0.5em]
\large Joint Emulation of Earth System Model Temperature--Precipitation Realizations with Internal Variability: fldgen v2.0\\[0.3em]
\normalsize OSTI ID: 1578031}
\author{Auto-generated Replication Blueprint}
\date{\today}

\begin{document}
\maketitle
\tableofcontents
\newpage

%---------------------------------------------------------------------
\section{Paper Overview}
%---------------------------------------------------------------------
This paper describes \texttt{fldgen} v2.0, an R package for generating spatially and temporally correlated realizations of temperature and precipitation fields that emulate the output of Earth System Models (ESMs). The method decomposes ESM output into a forced response (fitted via global-mean temperature scaling patterns) and residual variability, then generates new variability realizations preserving space--time correlation structure and cross-variable (temperature--precipitation) correlations using empirical orthogonal functions (EOFs) and a vector autoregressive (VAR) model.

%---------------------------------------------------------------------
\section{Detailed Workflow}
%---------------------------------------------------------------------

\subsection{Phase 1: Package Installation and Data Acquisition}
\begin{enumerate}[label=\textbf{1.\arabic*}]
  \item \textbf{Install fldgen v2.0} from GitHub.
  \item \textbf{Download ESM training data.}
    \begin{itemize}
      \item CMIP5/CMIP6 model output (temperature and precipitation fields).
      \item Global-mean temperature trajectories from simple climate models (e.g., Hector, MAGICC).
    \end{itemize}
  \item \textbf{Prepare input data} in the required NetCDF format.
\end{enumerate}

\subsection{Phase 2: Pattern Scaling (Forced Response)}
\begin{enumerate}[label=\textbf{2.\arabic*}]
  \item \textbf{Fit scaling patterns.}
    \begin{itemize}
      \item Regress gridded temperature and precipitation fields against global-mean temperature.
      \item Obtain per-grid-cell scaling coefficients (pattern fields).
    \end{itemize}
  \item \textbf{Compute residuals.}
    \begin{itemize}
      \item Subtract the scaled forced response from the ESM output to obtain variability residuals.
    \end{itemize}
\end{enumerate}

\subsection{Phase 3: Variability Model (EOF + VAR)}
\begin{enumerate}[label=\textbf{3.\arabic*}]
  \item \textbf{Joint EOF decomposition.}
    \begin{itemize}
      \item Concatenate temperature and precipitation residual fields.
      \item Compute EOFs (PCA) to reduce dimensionality; retain leading $K$ modes.
    \end{itemize}
  \item \textbf{Fit VAR model} to the EOF time series to capture temporal autocorrelation.
  \item \textbf{Generate new realizations.}
    \begin{itemize}
      \item Draw innovations from the fitted VAR noise distribution.
      \item Propagate through the VAR model to produce new EOF coefficient time series.
      \item Project back to physical space via EOF patterns.
      \item Add forced response to obtain full temperature/precipitation fields.
    \end{itemize}
\end{enumerate}

\subsection{Phase 4: Validation}
\begin{enumerate}[label=\textbf{4.\arabic*}]
  \item \textbf{Compare emulated statistics to ESM output.}
    \begin{itemize}
      \item Spatial correlation structure (variograms).
      \item Temporal autocorrelation functions.
      \item Cross-variable (T--P) correlations.
      \item Marginal distributions at selected grid cells.
    \end{itemize}
  \item \textbf{Reproduce paper figures.}
\end{enumerate}

%---------------------------------------------------------------------
\section{Datasets}
%---------------------------------------------------------------------

\begin{longtable}{p{4.5cm} p{3.5cm} p{6cm}}
\toprule
\textbf{Dataset / Source} & \textbf{Type} & \textbf{Location / Notes} \\
\midrule
\endfirsthead
\toprule
\textbf{Dataset / Source} & \textbf{Type} & \textbf{Location / Notes} \\
\midrule
\endhead
fldgen test/example data & Bundled & Included in the fldgen package \newline \url{https://github.com/jgcri/fldgen} \\
\midrule
CMIP5/CMIP6 ESM output & Public & \url{https://esgf-node.llnl.gov/search/cmip6/} \\
\midrule
Hector global-mean temperature & Public & \url{https://github.com/JGCRI/hector} \\
\midrule
MAGICC temperature pathways & Public & \url{https://magicc.org/} \\
\bottomrule
\end{longtable}

%---------------------------------------------------------------------
\section{Models, Tools, and Software}
%---------------------------------------------------------------------

\subsection{Programming Languages}
\begin{itemize}
  \item \textbf{R 4.0+} --- fldgen is an R package.
\end{itemize}

\subsection{Libraries and Frameworks}
\begin{longtable}{p{4.5cm} p{2.5cm} p{6.5cm}}
\toprule
\textbf{Tool / Library} & \textbf{Version} & \textbf{Purpose / URL} \\
\midrule
\endfirsthead
\toprule
\textbf{Tool / Library} & \textbf{Version} & \textbf{Purpose / URL} \\
\midrule
\endhead
fldgen & v2.0 & Core field generation package \newline \url{https://github.com/jgcri/fldgen} \\
\midrule
ncdf4 (R) & latest & NetCDF file I/O \\
\midrule
vars (R) & latest & Vector autoregressive model fitting \\
\midrule
raster / terra (R) & latest & Geospatial data handling \\
\midrule
ggplot2 (R) & latest & Visualization \\
\midrule
Hector & latest & Simple climate model for global-mean temperature \newline \url{https://github.com/JGCRI/hector} \\
\bottomrule
\end{longtable}

\subsection{Hardware Requirements}
\begin{itemize}
  \item Standard workstation; EOF computations on global grids benefit from $\geq$ 16\,GB RAM.
\end{itemize}

\end{document}
